\documentclass[11pt,a4paper]{article}
\usepackage[T1]{fontenc}
\usepackage[utf8]{inputenc}
\usepackage[margin=2.3cm]{geometry}
\usepackage{amsmath,amssymb}
\usepackage{booktabs}
\usepackage{microtype}
\usepackage{hyperref}
\usepackage{xcolor}
\hypersetup{colorlinks=true,linkcolor=black,urlcolor=blue!60!black}

\title{Probabilistic Day-Ahead Electricity Price Forecasting\\
       for Germany--Luxembourg and Spain}
\author{Team: Higher Power\\ ETH Z\"urich Frigg Hackathon}
\date{May 2026}

\newcommand{\DELU}{\textsf{DE-LU}}
\newcommand{\ES}{\textsf{ES}}

\begin{document}
\maketitle

\begin{abstract}
\noindent
We forecast hourly day-ahead electricity prices for the German--Luxembourg
(\DELU) and Spanish (\ES) bidding zones with calibrated probabilistic
intervals. The pipeline combines a gradient-boosted quantile regression
(short-term, $\le 14$ days) with an STL-based seasonal extrapolation
(long-term, $> 14$ days), trained on price, generation, weather, and fuel
data from 2024-01-01 onwards. Forecasts are scored by pinball loss at
$q \in \{0.025, 0.45, 0.975\}$ for each zone.
\end{abstract}

\section{Problem statement}
\label{sec:problem}

Let $p_t^z \in \mathbb{R}$ denote the day-ahead auction clearing price
(EUR/MWh) for hour $t$ in zone $z \in \{\DELU, \ES\}$. Given information
$\mathcal{F}_{t_0}$ available at decision time $t_0$, we produce, for
each hour $t \in [t_0+1\,\mathrm{d},\ t_0+2\,\mathrm{d}]$ in the
evaluation window, three quantile forecasts
$\hat{q}_{0.025}^z(t),\ \hat{q}_{0.45}^z(t),\ \hat{q}_{0.975}^z(t)$.

Forecast quality is scored by the pinball loss
\begin{equation}
  \rho_q(y,\hat{q})
  \;=\; \max\bigl\{q\,(y - \hat{q}),\; (q-1)(y - \hat{q})\bigr\},
  \qquad q \in (0,1),
  \label{eq:pinball}
\end{equation}
evaluated at $q = 0.45$. The slight asymmetry $q < 0.5$ penalises
\emph{over}-prediction more heavily than under-prediction, consistent
with a price-taker procuring electricity at the auction.

\section{Data}
\label{sec:data}

All inputs are public, hourly, and aligned to UTC.

\begin{table}[h]
\centering
\small
\begin{tabular}{@{}llll@{}}
\toprule
Source & Variables & Range & Notes \\ \midrule
Energy-Charts (Fraunhofer ISE) &
  $p_t^z$, generation per technology & 2024-01 -- live & DE-LU and ES \\
Open-Meteo Archive &
  $T_t$, $v_t$, $I_t$ & 2024-01 -- today & temp, wind, irradiance \\
Open-Meteo Forecast &
  $T_t$, $v_t$, $I_t$ & next $16$ days & ECMWF + DWD blend \\
Yahoo Finance &
  NG=F, KRBN ETF & daily, ffilled & gas / CO$_2$ proxies \\ \bottomrule
\end{tabular}
\caption{Data sources. Generation totals are decomposed into
$W_t$ (wind), $S_t$ (solar PV), $L_t$ (load), $G_t$ (fossil gas).}
\end{table}

Prices are clipped to $[-500, 3000]$ EUR/MWh and a 1-hour grid is
imposed by linear resampling. Up to three consecutive missing values
are forward-filled; longer gaps are imputed by column median.

\section{Feature engineering}
\label{sec:features}

For each zone and hour we construct $43$ features. With
$h(t),\ d(t),\ m(t)$ denoting hour-of-day, day-of-week, and month, all
periodic encodings use the standard sine/cosine pair
\begin{equation}
  \phi_P(t) = \bigl(\sin(2\pi\,k(t)/P),\ \cos(2\pi\,k(t)/P)\bigr),
\end{equation}
for periods $P \in \{24, 7, 12\}$.

\paragraph{Price autoregression.}
Lags of the target,
$\bigl\{\,p_{t-\tau}^z : \tau \in \{24, 48, 72, 120, 168\}\,\bigr\}$,
plus a one-week momentum $p_{t-24}^z - p_{t-168}^z$ and rolling
statistics over windows $W \in \{24, 48, 72, 168\}$,
\begin{equation}
  \mu_W(t) = \tfrac{1}{W}\!\!\sum_{i=1}^{W} p_{t-i}^z,
  \qquad
  \sigma_W(t) = \sqrt{\tfrac{1}{W-1}\!\!\sum_{i=1}^{W}(p_{t-i}^z - \mu_W(t))^2},
\end{equation}
together with $\min$ and the empirical $0.1$-quantile over $W=24,48$ to
capture price-floor behaviour.

\paragraph{Cross-zone leakage.}
DE-LU and ES are weakly coupled through their common weather macro-pattern
and via cross-border interconnectors (FR transit). We include
$p_{t-24}^{\bar{z}},\ p_{t-48}^{\bar{z}}$ where $\bar{z}$ is the partner
zone, which contributes $\sim$$3\%$ of LightGBM gain on \DELU.

\paragraph{Exogenous.}
Weather $(T_t, v_t, I_t)$, generation $(W_t, S_t, L_t, G_t)$, and
\begin{equation}
  \sigma_t \;=\; \frac{S_t}{\max(L_t,\ 1)},
  \qquad
  \Delta I_t \;=\; I_t - I_{t-1},
\end{equation}
the \emph{solar surplus ratio} and irradiance change. Fuel proxies
enter as 1- and 7-day-lagged daily-resampled values; binary holiday
flags for DE and ES come from the \texttt{holidays} package.

\section{Short-term model ($\le 14$ days)}
\label{sec:short}

For each zone $z$ and quantile $q \in \{0.025, 0.45, 0.975\}$ we fit a
gradient-boosted regressor $f_{\theta}^{z,q}: \mathbb{R}^{43} \to \mathbb{R}$
minimising the empirical pinball loss
\begin{equation}
  \mathcal{L}_q(\theta) \;=\;
  \sum_{t} w_t \,\rho_q\!\bigl(p_t^z,\ f_{\theta}^{z,q}(x_t)\bigr).
  \label{eq:loss}
\end{equation}
We use LightGBM with \texttt{objective='quantile'}, $\texttt{alpha}=q$,
$800$ trees (early-stopped), \texttt{learning\_rate} $=0.04$,
\texttt{num\_leaves} $=63$, $L_1=0.1$, $L_2=0.5$, and column/row
subsampling at $0.8$.

\subsection{Time-decayed sample weights}

Recent regimes are more informative than the 2022 gas crisis, so
samples are exponentially down-weighted with half-life
$H = 180$~days:
\begin{equation}
  w_t \;=\; \frac{\tilde{w}_t}{\bar{\tilde{w}}}, \qquad
  \tilde{w}_t \;=\; \exp\!\left(-\tfrac{\ln 2}{H}\,\Delta_t\right), \qquad
  \Delta_t \;=\; \frac{t_{\max} - t}{86400 \, \mathrm{s}}.
\end{equation}
Six-month-old samples thus enter at half the leverage of today's;
two-year-old samples at $\sim$$8\%$.

\subsection{Validation and ensembling}

The last $8 \times 168 = 1344$ hours are held out for early stopping
with patience $60$. The selected number of rounds $K^\star$ is then
\emph{re-fit on the full training set}; this exploits the additional
$8$ weeks of recent data without re-tuning. We repeat the procedure
for $M = 5$ random seeds and average:
\begin{equation}
  \hat{q}^{\,z,q}_{\mathrm{ST}}(t)
  \;=\; \frac{1}{M}\sum_{m=1}^{M} f_{\theta_m}^{z,q}(x_t).
  \label{eq:ensemble}
\end{equation}
Total: $2 \text{ zones} \times 3 \text{ quantiles} \times 5 \text{ seeds} = 30$
LightGBM models per run.

\subsection{Constructing exogenous features for the eval window}

Lag features at $t$ in the eval window simply read from history.
Generation features $(W_t, S_t, L_t, G_t)$ must be \emph{forecast}.
We use a month$\times$hour climatological mean
$\bar{X}_{m,h} = \mathbb{E}[X_t \mid m(t)=m,\ h(t)=h]$ as a baseline,
and rescale wind/solar by the realized weather-forecast ratio
\begin{equation}
  \hat{S}_t \;=\; \bar{S}_{m(t),h(t)} \cdot
  \mathrm{clip}\!\left(\frac{I_t}{\bar{I}_{m(t),h(t)}},\ 0,\ 3\right),
  \qquad
  \hat{W}_t \;=\; \bar{W}_{m(t),h(t)} \cdot
  \mathrm{clip}\!\left(\frac{v_t}{\bar{v}_{m(t),h(t)}},\ 0,\ 3\right),
\end{equation}
preserving variance from the meteorological forecast while anchoring
levels to seasonal climatology. Renewable share is then recomputed
from the adjusted $\hat{S} + \hat{W}$ over the (climatological) load.

\section{Long-term model ($> 14$ days)}
\label{sec:long}

Beyond two weeks the lag-driven short-term model has no signal:
$p_{t-168}$ becomes irrelevant. We use seasonal-trend decomposition
(STL, robust, weekly period $P = 168$):
\begin{equation}
  p_t^z \;=\; T_t + S_t + R_t,
\end{equation}
fitted on the full training history per zone.

\paragraph{Trend extrapolation.}
The trend is linearly extrapolated from the last $4$ weeks
($N = 28 \times 24 = 672$ hours):
\begin{equation}
  \hat{T}_{t_0+h} \;=\; T_{t_0} + \alpha\, h,
  \qquad
  \alpha \;=\; \arg\min_{\beta}
  \sum_{i=t_0-N+1}^{t_0} \bigl(T_i - (a + \beta \, i)\bigr)^2.
\end{equation}

\paragraph{Seasonal anchor.}
We collapse the STL seasonal to a month$\times$hour climatology
$\bar{S}_{m,h}$ and read off
$\hat{S}_{t} = \bar{S}_{m(t),\,h(t)}$.

\paragraph{Probabilistic intervals.}
Empirical $2.5\%$ and $97.5\%$ residual quantiles $r_{0.025}, r_{0.975}$
are widened with horizon (in days, $h_d$):
\begin{equation}
  \hat{q}^{\,z,p}_{\mathrm{LT}}(t)
  \;=\; \hat{T}_{t} + \hat{S}_t \;+\; r_p \cdot
  \sqrt{\max(h_d/14,\ 1)},
  \qquad p \in \{0.025, 0.975\},
\end{equation}
with $\hat{q}_{0.45}^{\mathrm{LT}}(t) = \hat{T}_t + \hat{S}_t$. The
$\sqrt{\cdot}$ scaling is the diffusive widening expected from a random-walk
residual.

\section{Composition and post-processing}

\subsection{Regime switching}

Let $h_d(t)$ denote the lead time in days from $t_0$. The combined
forecast is
\begin{equation}
  \hat{q}^{\,z,p}(t) \;=\;
  \begin{cases}
    \hat{q}^{\,z,p}_{\mathrm{ST}}(t) & \text{if } h_d(t) \le 14, \\
    \hat{q}^{\,z,p}_{\mathrm{LT}}(t) & \text{if } h_d(t) > 14.
  \end{cases}
\end{equation}
For the May~11 evaluation window all hours fall in the short-term
regime; the long-term branch covers extended-horizon submissions.

\subsection{Solar-driven interval widening}

Backtest analysis showed the LightGBM model under-estimates uncertainty
in high-irradiance midday hours, when negative-price events caused by
photovoltaic surplus dominate. We multiply the half-widths by
\begin{equation}
  \phi(I_t) \;=\;
  \min\!\left(1 + \frac{\max(I_t - 150,\ 0)}{400},\ 2.0\right),
\end{equation}
a piecewise-linear ramp from $1$ at $I = 150~\mathrm{W/m^2}$
to a cap of $2$ at $\sim 550~\mathrm{W/m^2}$, applied as
\begin{align}
  \hat{q}^{\,z,0.025\,\prime}(t)
    &= \hat{q}^{\,z,0.45}(t) - \phi(I_t)\bigl(\hat{q}^{\,z,0.45}(t) - \hat{q}^{\,z,0.025}(t)\bigr), \\
  \hat{q}^{\,z,0.975\,\prime}(t)
    &= \hat{q}^{\,z,0.45}(t) + \phi(I_t)\bigl(\hat{q}^{\,z,0.975}(t) - \hat{q}^{\,z,0.45}(t)\bigr).
\end{align}

\subsection{Quantile sorting}

Independent fitting per quantile can produce crossings
$\hat{q}_{0.025} > \hat{q}_{0.45}$. We enforce monotonicity by sorting
the triplet:
\begin{equation}
  (\hat{q}_{0.025}',\ \hat{q}_{0.45}',\ \hat{q}_{0.975}')
  \;:=\; \mathrm{sort}(\hat{q}_{0.025},\ \hat{q}_{0.45},\ \hat{q}_{0.975}).
\end{equation}

\section{Limitations and future work}

\paragraph{Day-of-week recalibration.}
The pinball asymmetry interacts poorly with weekend/holiday regime
shifts. A targeted recalibration --- e.g.\ a day-of-week bias
correction
$\hat{q}^{\prime}_{p}(t) = \hat{q}_p(t) - b^{z}_{d(t)}$ where
$b^{z}_{d(t)}$ is the median residual on matching weekday/weekend
samples --- would improve calibration across regimes.

\paragraph{Coverage-driven widening.}
Replacing the manual solar widening with an isotonic/conformal
recalibration on a holdout set would yield empirically valid
$95\%$ intervals across all regimes, not just sunny ones.

\paragraph{Joint zone modelling.}
Treating $(\,p^{\DELU}_t,\ p^{\ES}_t\,)$ as a bivariate target with a
shared latent renewables process should improve both bias and
interval calibration during European-wide weather events. A
multi-output LightGBM with chained quantile heads is a low-effort
upgrade.

\paragraph{Reproducibility.}
All data is fetched from public APIs and cached locally; the run is
seed-deterministic given fixed cache contents. End-to-end pipeline:
$\sim$~4~min on a 2024 MacBook (data cached) for the full $30$-model
ensemble.

\end{document}
